% ch08-workplan.tex – Chapter 8: work plan – WBS, RACI, Gantt (final round, Writer). RACI rows generated from data/wbs.csv.
\ChapterMission{המשימה: לפרק את הפרויקט לחבילות עבודה, להגדיר מי אחראי לכל משימה ולפרוס את הכול על ציר הזמן}
\ChapterIcon{info}
\chapter{תוכנית עבודה: פירוק, אחריות וגאנט}\label{chap:workplan}
\begin{scenebox}
בחדר הישיבות של מחלקת התנועה תלוי לוח ארוך עם פסים צבעוניים. מהנדסת מזיזה פס אחד שבוע ימינה,
ושלושה פסים אחרים נגררים אחריו. מישהו שואל מי חותם על הבדיקה בצומת. כולם מביטים בעמודה אחת.
\end{scenebox}

השאלה מי חותם היא העיקר בתוכנית עבודה: לכל משימה צריך להיות אחראי אחד ובעל סמכות אחד.
הפרק מפרק את הפרויקט לחמש חבילות עבודה, מגדיר תפקידים ואחריות, ופורס את המשימות על פני שנה, עם אבני דרך שבהן מתקבלות החלטות.

\section{מבנה פירוק העבודה}
מבנה פירוק העבודה (\en{Work Breakdown Structure, WBS}) באיור~\ref{fig:t3} מחלק את הפרויקט לחמש חבילות: ייזום ואישורים, חישה ונתונים, בקר ותצוגה, קמפיין שליטה חברתית, ופיילוט והערכה.
החלוקה עוקבת אחר ארכיטקטורת המערכת (פרק~\ref{chap:architecture}), כך שלכל שכבה יש חבילה ובעל תפקיד מוביל, וחבילת הייזום מטפלת באישור הרגולטורי, שהוא הסיכון הגדול ביותר (פרק~\ref{chap:risks}).

\begin{figure}[ht]
  \centering
  \inputfigure{figures/tikz/t3-wbs}
  \caption{מבנה פירוק העבודה: חמש חבילות עבודה ומשימותיהן}
  \label{fig:t3}
\end{figure}

\section{לוח זמנים ואבני דרך}
תרשים הגאנט באיור~\ref{fig:gantt} פורס את המשימות על פני שנים-עשר חודשים. שלושת החודשים הראשונים מוקדשים לבחירת אתר, למיפוי התקן ולתכנון, ובמקביל מתחילות גם בדיקות החיישנים; ההתקנה, האימות ובדיקות הקבלה נמשכות עד החודש השמיני; הפיילוט והקמפיין פועלים יחד עד החודש האחד-עשר, והחודש האחרון מוקדש לניתוח.

\begin{figure}[ht]
  \centering
  \inputfigure{figures/tikz/t5-gantt}
  \caption{תרשים גאנט של הפרויקט לאורך שנים-עשר חודשים, עם אבני הדרך}
  \label{fig:gantt}
\end{figure}

שלוש אבני דרך הן נקודות החלטה, והן מסומנות בתרשים:
\en{M1}, בסוף החודש החמישי: התקבל היתר להפעלת האדום המסתגל. אם ההיתר מתעכב, הפרויקט ממשיך עם הספירה בפינוי ועם הקמפיין בלבד, כלומר עם חלופה \en{A3}.
\en{M2}, בסוף החודש השמיני: בדיקות הקבלה עברו, והפיילוט מתחיל.
\en{M3}, בסוף החודש השנים-עשר: דוח ההערכה, הכולל את בדיקת ההשערות של פרק~\ref{chap:problem}, מוגש לעירייה.

\section{מטריצת אחריות}
טבלה~\ref{tab:raci} מגדירה לכל משימה מי מבצע (\en{R}), מי בעל הסמכות ונושא באחריות הסופית (\en{A}), עם מי מתייעצים (\en{C}) ומי מעודכן (\en{I}).
בכל שורה יש בדיוק מבצע אחד ובעל סמכות אחד. מנהל הפרויקט הוא בעל הסמכות ברוב המשימות, אך נציג העירייה הוא בעל הסמכות בכל מה שנוגע לצומת עצמו: האישור הרגולטורי, בדיקות הקבלה, הפעלת הקמפיין והדוח המסכם.

\begin{table}[ht]
  \centering
  \caption{מטריצת אחריות (\en{RACI}) לפי משימות. פרויקט = מנהל פרויקט; תנועה = מהנדס תנועה; חישה = מהנדס חישה ונתונים; עירייה = נציג העירייה; קמפיין = מנהל קמפיין}
  \label{tab:raci}
  \footnotesize
  \begin{tabularx}{\linewidth}{lXccccc}
    \toprule
    \headrow \thead{משימה} & \thead{תיאור} & \thead{פרויקט} & \thead{תנועה} & \thead{חישה} & \thead{עירייה} & \thead{קמפיין} \\
    \midrule
    \num{1.1} & סקר אתרים ונתוני בסיס & \en{A} & \en{R} & \en{I} & \en{C} & – \\
    \num{1.2} & מיפוי התקן והגבולות הבטוחים & \en{A} & \en{R} & \en{C} & \en{I} & – \\
    \num{1.3} & הגשה לאישור האדום המסתגל & \en{R} & \en{C} & – & \en{A} & \en{I} \\
    \midrule
    \num{2.1} & בחירת חישה ובדיקת מעבדה & \en{A} & \en{C} & \en{R} & \en{I} & – \\
    \num{2.2} & שכבת נתונים ופרטיות & \en{A} & – & \en{R} & \en{C} & \en{I} \\
    \num{2.3} & התקנת חיישנים וכיול אזורים & \en{I} & \en{A} & \en{R} & \en{C} & – \\
    \num{2.4} & אימות דיוק הזיהוי & \en{A} & \en{C} & \en{R} & \en{I} & – \\
    \midrule
    \num{3.1} & אלגוריתם האדום המסתגל & \en{A} & \en{R} & \en{C} & \en{I} & – \\
    \num{3.2} & סימולציה ועיכוב כלי רכב & \en{A} & \en{R} & \en{I} & \en{C} & – \\
    \num{3.3} & שילוב ספירה במופע הפינוי & \en{A} & \en{R} & \en{C} & \en{I} & – \\
    \num{3.4} & בדיקות קבלה בצומת & \en{I} & \en{R} & \en{C} & \en{A} & – \\
    \midrule
    \num{4.1} & מסרים ועיצוב & \en{A} & \en{I} & – & \en{C} & \en{R} \\
    \num{4.2} & הפעלת הקמפיין & \en{C} & – & \en{I} & \en{A} & \en{R} \\
    \num{4.3} & סקר מודעות & \en{A} & – & \en{C} & \en{I} & \en{R} \\
    \midrule
    \num{5.1} & מדידת מצב קיים & \en{A} & \en{C} & \en{R} & \en{I} & – \\
    \num{5.2} & הפעלת פיילוט & \en{A} & \en{R} & \en{C} & \en{I} & – \\
    \num{5.3} & ניתוח ודוח סיום & \en{R} & \en{C} & – & \en{A} & \en{I} \\
    \bottomrule
  \end{tabularx}
\end{table}
